% =============================================================================
%  lec07.tex —— 第 7 讲 极限与余极限：函子性（底本 PDF 第 58–65 页）
% =============================================================================
\lecture{7}{极限与余极限：函子性}
\begin{center}\small\color{black!55}（底本 PDF 第 58–65 页）\end{center}

上一讲定义了 $\infty$-范畴中的极限与余极限。本讲研究它们的函子性。

\section{对 $u$ 的函子性}

\begin{construction}\label{L7-1.1}
设 $K$ 是单纯集，$\cC$ 是 $\infty$-范畴。对 $\Fun(K,\cC)$ 中的态射 $\alpha\colon u_0\to u_1$，有典范函子
\[
\cC_{/u_0}\to\cC_{/u_1},
\]
它与到 $\cC$ 的遗忘函子相容（见第 6 讲命题—构造 3.15）。利用这个函子，我们把 $\lim u_0\in\cC$ 典范地提升到 $\cC_{/u_1}$。由终对象的定义，存在本质唯一的态射
\[
\lim u_0\to\lim u_1,
\]
它与二者在 $\cC_{/u_1}$ 中的提升相容。

对偶地，存在本质唯一的态射
\[
\colim u_0\to\colim u_1,
\]
它与二者在 $\cC_{u_1/}$ 中的提升相容。

由同前所引，若 $\alpha$ 可逆，则这些态射也可逆。
\end{construction}

\nitem\label{L7-1.2} 在以后的讲次中，我们会把上述构造升级为函子
\[
\lim\colon\Fun(K,\cC)'\to\cC,
\]
其中 $\Fun(K,\cC)'\subset\Fun(K,\cC)$ 是由使 $\lim u$ 存在的图 $u\colon K\to\cC$ 构成的全子 $\infty$-范畴。

\section{对 $\cC$ 的函子性}

\begin{construction}\label{L7-2.1}
设 $F\colon\cC\to\cC'$ 是 $\infty$-范畴之间的函子，$u\colon K\to\cC$ 是图。有显然的函子
\[
\cC_{/u}\to\cC'_{/F\circ u},
\]
它借助遗忘函子与 $F$ 相容。利用这个函子，我们把 $F(\lim u)\in\cC'$ 典范地提升到 $\cC'_{/F\circ u}$。由终对象的定义，存在本质唯一的态射
\[
F(\lim u)\to\lim(F\circ u),
\]
它与二者在 $\cC'_{/F\circ u}$ 中的提升相容。

对偶地，存在本质唯一的态射
\[
\colim(F\circ u)\to F(\colim u),
\]
它与二者在 $\cC'_{F\circ u/}$ 中的提升相容。
\end{construction}

\begin{definition}\label{L7-2.2}
设 $F\colon\cC\to\cC'$ 是 $\infty$-范畴之间的函子。若极限图 $\overline u\colon K^{\lhd}\to\cC$ 的复合 $F\circ\overline u\colon K^{\lhd}\to\cC'$ 仍是极限图，则称该极限图被 $F$ \textbf{保持}（preserve）；此时也说 $\lim u$ 被 $F$ 保持。
\yizhu{原书用 $K^{\lhd}$ 记在 $K$ 的末端接上一个终对象所得的锥（即扩张图 $K^{\lhd}\to\cC$ 以新加的点为锥顶）。不少文献把这种“右锥”记作 $K^{\triangleright}$，而把 $\Delta^0\star K$ 记作 $K^{\lhd}$；译文照原书记号，不改。\textbf{[来源：B]}}
\end{definition}

\begin{exercise}\label{L7-2.3}
证明：$F$ 保持 $\lim u$ 当且仅当 $F(\lim u)\to\lim(F\circ u)$ 是同构。
\end{exercise}

\begin{exercise}\label{L7-2.4}
证明：若两个函子 $F,G\colon\cC\to\cC'$ 彼此等价，则一个极限被 $F$ 保持当且仅当被 $G$ 保持。
\end{exercise}

\begin{exercise}\label{L7-2.5}
可逆函子保持所有存在的极限。
\end{exercise}

\nitem\label{L7-2.6} 尽管我们尚未引入 $\infty$-范畴之间的伴随函子，还是先记下下述结果：

\begin{theorem}\label{L7-2.7}
设 $F\colon\cC\rightleftarrows\cD\colon G$ 是 $\infty$-范畴之间的伴随。则
\begin{enumerate}[label=(\arabic*),leftmargin=2.4em,itemsep=0.1em]
\item 函子 $F$ 保持所有存在的余极限。
\item 函子 $G$ 保持所有存在的极限。
\end{enumerate}
\end{theorem}

\begin{definition}\label{L7-2.8}
设 $F_i\colon\cC\to\cD_i$ 是一族 $\infty$-范畴之间的函子。对图 $u\colon K\to\cC$，若只要 $F_i\circ\overline u\colon K^{\lhd}\to\cD_i$ 都是极限图，就有扩张图 $\overline u\colon K^{\lhd}\to\cC$ 是极限图，则称 $\{F_i\}$ \textbf{检测}（detect）$u$ 的极限。
\end{definition}

\begin{warning}\label{L7-2.9}
上述定义中，我们并不先验地假定 $\lim u$ 存在。
\end{warning}

\nitem\label{L7-2.10} 尽管我们尚未定义可表函子，还是先记下下述结果：

\begin{theorem}\label{L7-2.11}
余可表函子保持并检测所有极限。
\end{theorem}

\begin{remark}\label{L7-2.12}
下面把该定理说得更具体些。在以后的讲次中，我们会构造典范函子\footnote{我们在第 4 讲第 10 节只构造了函子 $\cC^{\op}\times\cC\to\hGrpd_\infty$。}
\[
\Map(-,-)\colon\cC^{\op}\times\cC\to\Grpd_\infty.
\]
于是该定理说的是：
\begin{enumerate}[label=(\arabic*),leftmargin=2.4em,itemsep=0.1em]
\item 图 $\overline u\colon K^{\lhd}\to\cC$ 是极限图，当且仅当对任意对象 $x\in\cC$，协变函子 $\Map(x,-)$ 把它送到 $\Grpd_\infty$ 中的极限图。
\item 图 $\overline u\colon K^{\lhd}\to\cC^{\op}$ 是极限图，当且仅当对任意对象 $x\in\cC$，反变函子 $\Map(-,x)$ 把它送到 $\Grpd_\infty$ 中的极限图。
\end{enumerate}
\end{remark}

\begin{warning}\label{L7-2.13}
若把 $\Grpd_\infty$ 换成 $\hGrpd_\infty$，类似的断言会不成立。
\end{warning}

\begin{theorem}\label{L7-2.14}
（Ker.02X9）设 $\cC$ 是 $\infty$-范畴，$B$ 是任意单纯集。则 $\Fun(B,\cC)$ 中的（余）极限可以逐点计算。
\end{theorem}

\begin{remark}\label{L7-2.15}
上述定理的确切含义如下。设 $u\colon K\to\Fun(B,\cC)$ 是任意图。若对 $B$ 中任意 $0$-单形 $b$，图
\[
u|_b\colon K\to\Fun(B,\cC)\xrightarrow{ev_b}\cC
\]
有极限，则 $\lim u$ 存在，且被函子族 $ev_b\colon\Fun(B,\cC)\to\cC$ 保持并检测。特别地，有
\[
(\lim u)|_b\xrightarrow{\ \simeq\ }\lim u|_b.
\]
注意这个断言并非平凡，因为函子 $B\to\cC$，$b\mapsto\lim u|_b$ 的构造并不平凡。见第 1 节。
\end{remark}

\begin{proposition}\label{L7-2.16}
全忠实函子检测所有（余）极限。
\end{proposition}

\begin{exercise}\label{L7-2.17}
上述命题的证明以下述练习的形式给出梗概。设 $F\colon\cC\to\cC'$ 是函子。证明：
\begin{enumerate}[label=(\arabic*),leftmargin=2.4em,itemsep=0.1em]
\item 函子 $F$ 全忠实，当且仅当它可由拟范畴之间的单态射 $F\colon\cC\to\cC'$ 表示，并且 $\cC'$ 的一个 $n$-单形落在 $F$ 的像中，当且仅当该单形的所有顶点都具有同一性质（即落在 $F$ 的像中）。
\item 若 $F$ 全忠实，则对任意单纯集 $K$，$\Fun(K,\cC)\to\Fun(K,\cC')$ 全忠实。
\item 若 $F$ 全忠实，则对任意图 $u\colon K\to\cC$，$\cC_{/u}\to\cC'_{/F\circ u}$ 全忠实。
\item 若 $F$ 全忠实，则它检测极限。
\end{enumerate}
\end{exercise}

\begin{warning}\label{L7-2.18}
全忠实函子未必保持极限。例如：$[0]\xrightarrow{0}[1]$ 不保持空极限，即终对象。
\end{warning}

\section{对 $K$ 的函子性}

\begin{construction}\label{L7-3.1}
设 $v\colon K'\to K$ 是 $\Set_\Delta$ 中的态射，$u\colon K\to\cC$ 是 $\infty$-范畴 $\cC$ 中的图。有限制函子
\[
\cC_{/u}\to\cC_{/u\circ v},
\]
它与到 $\cC$ 的遗忘函子相容。利用这个函子，我们把 $\lim_K u\in\cC$ 典范地提升到 $\cC_{/u\circ v}$。由终对象的定义，存在本质唯一的态射
\[
\lim_K u\to\lim_{K'}(u\circ v),
\]
它与二者在 $\cC_{/u\circ v}$ 中的提升相容。

对偶地，存在本质唯一的态射
\[
\colim_{K'}(u\circ v)\to\colim_K u,
\]
它与二者在 $\cC_{u\circ v/}$ 中的提升相容。
\end{construction}

\begin{definition}\label{L7-3.2}
设 $v\colon K'\to K$ 是 $\Set_\Delta$ 中的态射。若对任意 $\infty$-范畴 $\cC$ 与任意图 $u\colon K\to\cC$，限制函子 $\cC_{/u}\to\cC_{/u\circ v}$ 都是 $\infty$-范畴之间的等价\footnote{由 Ker.02NR，这个定义与 Joyal 和 Lurie 著作中的定义等价（见 Ker.02N1）。}，则称 $v$ 是\textbf{初始的}\footnote{文献中的其他术语：op-cofinal、left cofinal。}。

对偶地，若对任意 $\infty$-范畴 $\cC$ 与任意图 $u\colon K\to\cC$，限制函子 $\cC_{u/}\to\cC_{u\circ v/}$ 都是 $\infty$-范畴之间的等价，则称 $v\colon K'\to K$ 是\textbf{终结的}\footnote{文献中的其他术语：cofinal、right cofinal。}。
\end{definition}

\nitem\label{L7-3.3} 显然，若 $v$ 初始，则对任意图 $u\colon K\to\cC$，$\lim_K u\to\lim_{K'}(u\circ v)$ 可逆。事实上，逆命题也成立（命题~\ref{L7-3.12}）。

\nitem\label{L7-3.4} 注意 $v\colon K'\to K$ 初始，当且仅当 $v^{\op}\colon(K')^{\op}\to K^{\op}$ 终结。因此下面我们只讨论初始态射。

\begin{proposition}\label{L7-3.5}
（Ker.02NN，02NP）对 $\Set_\Delta$ 中的态射，“初始”这一性质在等价下不变。
\end{proposition}

\nitem\label{L7-3.6} 由此，我们得到 $\infty$-范畴之间初始函子的概念。

\begin{warning}\label{L7-3.7}
普通范畴之间的初始函子\footnote{关于其含义，见 [Sta24, Tag 09WN]。}未必是 $\infty$-范畴之间的初始函子。
\end{warning}

\begin{proposition}\label{L7-3.8}
（Ker.043E）拟范畴之间的函子 $x\colon\Delta^0\to K$ 初始，当且仅当 $x$ 是 $K$ 中的始对象。
\end{proposition}

\begin{warning}\label{L7-3.9}
一般地，
\begin{itemize}[leftmargin=2.2em,itemsep=0.1em]
\item 初始态射 $v\colon K'\to K$ 未必是 $\Fun(K',K)$ 中的始对象。例如：$\id\colon\Delta^1\to\Delta^1$。
\item $\Fun(K',K)$ 中的始对象未必是初始态射 $v\colon K'\to K$。例如：$K'\to\Delta^0$，其中 $K'$ 不是弱可缩的。见下面的命题~\ref{L7-3.19}。
\end{itemize}
\end{warning}

\nitem\label{L7-3.10} 作为一个应用，我们得到下述有用的极限判据。

\begin{proposition}\label{L7-3.11}
设 $\overline u\colon K^{\lhd}\to\cC$ 是 $\infty$-范畴 $\cC$ 中的图。记 $u\coloneqq\overline u|_K$。考虑限制函子
\[
\cC_{/\overline u(*)}\leftarrow\cC_{/\overline u}\to\cC_{/u}.
\]
\begin{enumerate}[label=(\arabic*),leftmargin=2.4em,itemsep=0.1em]
\item 函子 $\cC_{/\overline u(*)}\leftarrow\cC_{/\overline u}$ 是等价。
\item 函子 $\cC_{/\overline u}\to\cC_{/u}$ 是等价，当且仅当 $\overline u$ 是极限图。
\end{enumerate}
特别地，只要 $\lim u$ 存在，就得到等价 $\cC_{/\lim u}\xrightarrow{\ \simeq\ }\cC_{/u}$。
\end{proposition}

\begin{proof}
（1）由锥顶 $\Delta^0\to K^{\lhd}$ 初始（命题~\ref{L7-3.8}）可得。只需再证（2）。由第 6 讲练习 3.8，有典范等价 $\cC_{/\overline u}\simeq(\cC_{/u})_{/\overline u(*)}$。于是断言由第 6 讲命题 4.2 推出。
\end{proof}

\begin{proposition}\label{L7-3.12}
设 $v\colon K'\to K$ 是 $\Set_\Delta$ 中的态射。下列条件等价：
\begin{enumerate}[label=(\arabic*),leftmargin=2.4em,itemsep=0.1em]
\item 态射 $v$ 是初始的。
\item 对 $\infty$-范畴 $\cC$ 中的任意图 $u\colon K\to\cC$，都有：
\begin{itemize}[leftmargin=2.2em,itemsep=0.1em]
\item $u$ 的极限存在，当且仅当 $u\circ v$ 的极限存在；
\item 态射 $\lim_K u\to\lim_{K'}(u\circ v)$ 可逆。
\end{itemize}
\end{enumerate}
\end{proposition}

\begin{proof}
（1）$\Rightarrow$（2）显然。现设 $v$ 满足（2）。

设 $\cC$ 是任意 $\infty$-范畴，要证 $\cC_{/u}\to\cC_{/u\circ v}$ 是等价。在以后的讲次中我们会证明，存在全忠实函子 $\cC\to\cD$，使得 $\cD$ 具有所有以 $K$ 为指标的极限\footnote{只要 $K$ 与 $\cC$ 都小，Yoneda 嵌入 $\cC\to\Fun(\cC,\Grpd_\infty)^{\op}$ 就是这样的一个选择。}。考虑复合 $w\colon K\xrightarrow{u}\cC\to\cD$。我们断言
\begin{itemize}[leftmargin=2.2em,itemsep=0.1em]
\item 若 $\cD_{/w}\to\cD_{/w\circ v}$ 可逆，则 $\cC_{/u}\to\cC_{/u\circ v}$ 也可逆。
\end{itemize}
事实上，在此假设下，有交换图
\begin{center}
\begin{tikzcd}[column sep=4.6em,row sep=2.8em]
\cC_{/u} \arrow[r] \arrow[d] & \cC_{/u\circ v} \arrow[d] \\
\cD_{/w} \arrow[r, "\simeq"] & \cD_{/w\circ v}
\end{tikzcd}
\end{center}
由练习~\ref{L7-2.17}(3)，竖直函子全忠实，且借助底部的等价容易识别它们的本质像。由此形式地推出，顶部函子如所愿是等价。

于是只需证明 $\cD_{/w}\to\cD_{/w\circ v}$ 可逆。换言之，可以假定 $\cC$ 具有所有以 $K$ 为指标的极限。利用命题~\ref{L7-3.11}，容易构造交换图
\begin{center}
\begin{tikzcd}[column sep=5em,row sep=3em]
\cC_{/\lim u} \arrow[r, "\simeq"] \arrow[d, "\simeq"] & \cC_{/\lim(u\circ v)} \arrow[d, "\simeq"] \\
\cC_{/u} \arrow[r] & \cC_{/u\circ v},
\end{tikzcd}
\end{center}
它蕴含所要的断言。
\end{proof}

\begin{exercise}\label{L7-3.13}
设 $I\xrightarrow{f}J\xrightarrow{g}K$ 是 $\Set_\Delta$ 中的态射。设 $f$ 初始。则 $g$ 初始，当且仅当 $g\circ f$ 初始。
\end{exercise}

\begin{warning}\label{L7-3.14}
初始态射的集合不具有三选二性质。
\end{warning}

\begin{exercise}\label{L7-3.15}
考虑 $\Set_\Delta$ 中的下图
\begin{center}
\begin{tikzcd}[column sep=3.4em,row sep=2.8em]
& \Delta^1 \arrow[dr] & \\
\Delta^0 \arrow[ur, "1"] \arrow[rr] & & \Delta^0.
\end{tikzcd}
\end{center}
证明：三个态射中恰有两个是初始的。
\end{exercise}

\begin{proposition}\label{L7-3.16}
（Ker.02NK）终结（相应地，初始）态射的集合对有限积封闭。
\end{proposition}

\begin{remark}\label{L7-3.17}
一旦有了极限的分配律，上述命题便可确立。分配律说：对图 $u\colon J\times K\to\cC$，
\[
\lim_{J\times K}u\xrightarrow{\ \simeq\ }\lim_{j\in J}\Big[\lim_K u(j,-)\Big].
\]
然而这个断言并非平凡，因为函子 $J\to\cC$，$j\mapsto\lim_K u(j,-)$ 的构造并不平凡。见第 1 节。
\end{remark}

\nitem\label{L7-3.18} 第 6 讲命题 2.10 与命题 2.12 有下述推广：

\begin{proposition}\label{L7-3.19}
（Ker.02N5）设 $v\colon K'\to K$ 是 $\Set_\Delta$ 中的态射。
\begin{enumerate}[label=(\arabic*),leftmargin=2.4em,itemsep=0.1em]
\item 若 $K$ 是 Kan 复形，则 $v$ 初始（相应地，终结），当且仅当它是弱同伦等价。
\item 若 $v$ 初始（相应地，终结），则 $v$ 是弱同伦等价。
\end{enumerate}
\end{proposition}

\begin{exercise}\label{L7-3.20}
设 $K$ 是弱可缩单纯集，$\cC$ 是 $\infty$-范畴。对对象 $x\in\cC$，考虑取值 $x$ 的常值图 $K\to\cC$。证明：
\[
\colim_K x\to x\to\lim_K x
\]
都是同构。
\end{exercise}

\begin{exercise}\label{L7-3.21}
证明：对 $\cC\coloneqq\Grpd_\infty$，上述练习的断言对一般的 $K$ 会不成立。
\end{exercise}

\begin{proposition}\label{L7-3.22}
（Ker.02N9）若函子 $F\colon\cC\to\cD$ 把 $\cD$ 实现为 $\cC$ 关于其中一族态射的 $\infty$-范畴局部化，则 $F$ 既初始又终结。
\end{proposition}

\begin{exercise}\label{L7-3.23}
设 $u\colon K\to\cC$ 是 $\infty$-范畴 $\cC$ 中的图，且其像包含在 $\cC^{\simeq}$ 中。设 $K$ 弱可缩，证明：对任意 $i\in K$，
\[
u(i)\to\colim_K u,\qquad \lim_K u\to u(i)
\]
都是同构。
\end{exercise}

\nitem\label{L7-3.24} 我们再把下述结果记录下来：

\begin{theorem}\label{L7-3.25}
设 $F\colon\cC\rightleftarrows\cD\colon G$ 是 $\infty$-范畴之间的伴随。则
\begin{enumerate}[label=(\arabic*),leftmargin=2.4em,itemsep=0.1em]
\item 函子 $F$ 初始。
\item 函子 $G$ 终结。
\end{enumerate}
\end{theorem}

\section{Quillen 定理 A}

下述结果称为 $\infty$-范畴的 Quillen 定理 A，最早由 Joyal 证明。我们将在以后的讲次中证明它。

\begin{theorem}\label{L7-4.1}
设 $v\colon K'\to K$ 是单纯集之间的态射。设 $K$ 是拟范畴，则：
\begin{enumerate}[label=(\arabic*),leftmargin=2.4em,itemsep=0.1em]
\item $v$ 初始，当且仅当对任意对象 $x\in K$，$\Set_\Delta$ 中的纤维积 $K'\times_K K_{/x}$ 弱可缩。
\item $v$ 终结，当且仅当对任意对象 $x\in K$，$\Set_\Delta$ 中的纤维积 $K'\times_K K_{x/}$ 弱可缩。
\end{enumerate}
\end{theorem}

\begin{remark}\label{L7-4.2}
当 $K'=\Delta^0$ 时，定理退化为命题~\ref{L7-3.8}。这源于如下事实：$\Hom^L_K(x,y)\coloneqq\{y\}\times_K K_{x/}$ 是 $\infty$-群胚 $\Map_K(x,y)$ 的一个模型。
\end{remark}

\begin{remark}\label{L7-4.3}
当 $K=\Delta^0$ 时，定理退化为命题~\ref{L7-3.19}(1)。
\end{remark}

\begin{remark}\label{L7-4.4}
当 $K'$ 也是拟范畴时，在普通范畴 $\Set_\Delta$ 中取的纤维积 $K'\times_K K_{/x}$ 也计算拟范畴 $\mathrm{QCat}$ 中相应的纤维积（见第 6 讲脚注 4）。于是得到上述定理的 $\infty$-范畴版本：对 $\infty$-范畴之间的函子 $v\colon K'\to K$，
\begin{enumerate}[label=(\arabic*),leftmargin=2.4em,itemsep=0.1em]
\item $v$ 初始，当且仅当对任意对象 $x\in K$，纤维积 $K'\times_K K_{/x}$ 弱可缩。
\item $v$ 终结，当且仅当对任意对象 $x\in K$，纤维积 $K'\times_K K_{x/}$ 弱可缩。
\end{enumerate}
\end{remark}

\begin{remark}\label{L7-4.5}
在以后的讲次中，我们将证明：
\begin{enumerate}[label=(\arabic*),leftmargin=2.4em,itemsep=0.1em]
\item 函子 $v\colon K'\to K$ 有右伴随，当且仅当对任意对象 $x\in K$，纤维积 $K'\times_K K_{/x}$ 有终对象。
\item 函子 $v\colon K'\to K$ 有左伴随，当且仅当对任意对象 $x\in K$，纤维积 $K'\times_K K_{x/}$ 有始对象。
\end{enumerate}
\end{remark}

于是定理~\ref{L7-3.25}可由定理~\ref{L7-4.1}与第 6 讲命题 2.12 推出。

\begin{exercise}\label{L7-4.6}
设 $\Delta_{\mathrm{inj}}$ 是 $\Delta$ 的如下子范畴：
\begin{itemize}[leftmargin=2.2em,itemsep=0.1em]
\item 对象是 $[n]$，$n\ge 0$；
\item 态射是单的不减映射 $[m]\to[n]$。
\end{itemize}
证明 $N_{\bullet}(\Delta_{\mathrm{inj}})\to N_{\bullet}(\Delta)$ 初始。
\end{exercise}

\appendixsection{所有函子保持的极限}

\nitem\label{L7-A.1} 某些极限被所有函子保持。

\begin{construction}\label{L7-A.2}
考虑如下定义的普通范畴 $\mathrm{Idem}$：
\begin{itemize}[leftmargin=2.2em,itemsep=0.1em]
\item 只有一个对象 $\ast$；
\item $\Hom(\ast,\ast)\coloneqq\{\id,e\}$，其中 $e\circ e=e$。
\end{itemize}
\end{construction}

\begin{exercise}\label{L7-A.3}
找出 $\mathrm{Idem}$ 的脉的全部非退化单形。
\end{exercise}

\begin{exercise}\label{L7-A.4}
证明下述定理：
\end{exercise}

\begin{theorem}\label{L7-A.5}
以 $\mathrm{Idem}$ 为指标的极限与余极限被 $\infty$-范畴之间的任意函子保持。
\end{theorem}

\begin{construction}\label{L7-A.6}
设 $\Delta_{\mathrm{aug}}$ 是如下定义的普通范畴：
\begin{itemize}[leftmargin=2.2em,itemsep=0.1em]
\item 对象是 $[n]$，$n\ge -1$，其中 $[-1]\coloneqq\varnothing$；
\item 态射是不减映射 $[m]\to[n]$。
\end{itemize}
设 $\Delta_{\mathrm{split}}$ 是如下定义的普通范畴：
\begin{itemize}[leftmargin=2.2em,itemsep=0.1em]
\item 对象是 $[n]\sqcup\{-\infty\}$，$n\ge -1$；
\item 态射是把 $-\infty$ 送到 $-\infty$ 的不减映射 $[m]\sqcup\{-\infty\}\to[n]\sqcup\{-\infty\}$。
\end{itemize}
有显然的函子
\[
\Delta\to\Delta_{\mathrm{aug}}\to\Delta_{\mathrm{split}}.
\]
\end{construction}

\begin{exercise}\label{L7-A.7}
证明 $\Delta_{\mathrm{aug}}$ 可与 $\Delta^{\lhd}$ 等同。
\end{exercise}

\begin{theorem}\label{L7-A.8}
证明：若图 $u\colon\Delta^{\lhd}\to\cC$ 能扩张为图 $u'\colon\Delta_{\mathrm{split}}\to\cC$，则 $u$ 是极限图。此时称该极限为\textbf{分裂余单纯极限}。
\end{theorem}

\begin{exercise}\label{L7-A.9}
推出：分裂余单纯极限被任意函子保持。
\end{exercise}

\nitem\label{L7-A.10} \textbf{建议阅读}\quad Ker.03Y9，HA.4.7.2。

\appendixsection{$\Grpd_\infty$ 中的常值（余）极限}

\begin{exercise}\label{L7-B.1}
设 $K\in\Set_\Delta$，$\underline X\colon K\to\Kan$ 是常值 $X\in\Kan$ 的图。你能否把它扩张为图 $K^{\lhd}\to\Kan$，使得锥顶的取值是 $\Fun(K,X)$？
\end{exercise}

\begin{exercise}\label{L7-B.2}
设 $K\in\Set_\Delta$，$X,Y\in\Kan$。考虑常值函子 $\underline X,\underline Y\in\Fun(K,\Kan)$。构造 $\Kan$ 复形之间的下述弱等价：
\begin{align*}
\Fun(Y,\Fun(K,X)) &\simeq \Fun(K,\Fun(Y,X)) \simeq \Fun(K,\Hom^{\mathrm{L}}_{\Kan}(Y,X)) \\
&\simeq \Fun(K,\Hom^{\mathrm{B}}_{\Kan}(Y,X)) \simeq \Hom^{\mathrm{B}}_{\Fun(K,\Kan)}(\underline Y,\underline X).
\end{align*}
由此推出
\[
\lim_K \underline X\simeq\Fun(K,X).
\]
\end{exercise}

\begin{exercise}\label{L7-B.3}
设 $K\in\Set_\Delta$，$\underline X\colon K\to\Kan$ 是取值 $X\in\Kan$ 的常值函子。设 $K$ 是 Kan 复形，证明
\[
\colim_K \underline X\simeq K\times X.
\]
\end{exercise}

\subsection*{参考文献}
\par\noindent\hangindent=5.2em\hangafter=1 [Sta24]\hspace{0.6em}The Stacks project authors. The stacks project. \url{https://stacks.math.columbia.edu}, 2024.
